\documentclass[11pt]{article}
\usepackage[a4paper,margin=1in]{geometry}
\usepackage{amsmath,amssymb,amsthm}
\usepackage{graphicx}
\usepackage[colorlinks=true,linkcolor=blue,citecolor=blue,urlcolor=blue]{hyperref}

\newtheorem{theorem}{Theorem}
\newtheorem{lemma}[theorem]{Lemma}
\newtheorem{remark}[theorem]{Remark}
\newcommand{\N}{\mathbb N}
\newcommand{\K}{K_{\mathrm{SK}}}
\DeclareMathOperator{\supp}{supp}

\title{Positive Strongly Kreiss Weighted Shifts Have\\
Polylogarithmic Power Growth}
\author{Run \texttt{fa\_banach\_001}, lane 0}
\date{19 July 2026}

\begin{document}
\maketitle

\begin{center}
\textbf{Status: partial result, likely valid; human review requested.}
\end{center}

\section{Source question}

Arnold and Cuny ask the following in Question 5.8 of
\cite{ArnoldCuny}: if $1<p<\infty$ and $T$ is a positive strongly Kreiss
bounded operator on an $L^p$-space, must there be $C,\kappa>0$ such that
\[
  \|T^N\|\le C\log^\kappa(N+1)\qquad(N\in\N)?
\]
The source question is reproduced below from page 19 of the arXiv PDF.

\begin{center}
\includegraphics[width=\textwidth]{figures/open_problem_crop.png}
\end{center}

We prove an affirmative answer for positive unilateral weighted backward
shifts whose weights are at least one.  This class includes the familiar
telescoping logarithmic examples discussed in \cite{ArnoldCuny,CCEL}.

\section{Definitions}

Let $1\le p\le\infty$ and use the convention $\N_0=\{0,1,2,\ldots\}$.
Given a bounded sequence $(w_j)_{j\ge0}$ of positive numbers, the unilateral
weighted backward shift $T$ on $\ell^p(\N_0)$ is
\[
  (Tx)_j=w_jx_{j+1}.
\]
Put $a_0=1$ and $a_n=\prod_{j=0}^{n-1}w_j$ for $n\ge1$.  Then
\begin{equation}\label{eq:powers}
  (T^m x)_j=\frac{a_{j+m}}{a_j}x_{j+m},\qquad
  \|T^m\|=\sup_{j\ge0}\frac{a_{j+m}}{a_j}.
\end{equation}

The strong Kreiss condition is the iterated resolvent estimate from
\cite{ArnoldCuny}.  We use its equivalent exponential formulation: there is
$\K\ge1$ such that
\begin{equation}\label{eq:exp}
  \|e^{zT}\|\le \K e^{|z|}\qquad(z\in\mathbb C).
\end{equation}
This equivalence is recalled in the introduction of \cite{ArnoldCuny}.

\section{Proof intuition}

For a positive shift, $e^{-N}e^{NT}$ is a Poisson average of its powers.
Poisson mass is of order $N^{-1/2}$ throughout a window of width of order
$\sqrt N$ around $N$.  Test this average on a constant vector supported on
one coordinate block of that width and inspect its image on a second block.
Monotonicity of the cumulative weights makes every matrix entry between the
two blocks at least a fixed Poisson factor times
$a_{n+N}/a_{n+c\sqrt N}$.  There are $c\sqrt N$ entries in every relevant
row, exactly cancelling the $N^{-1/2}$ Poisson scale in every $\ell^p$ norm.

The strong Kreiss bound therefore forces the cumulative weight to grow by
only a fixed factor when its index is reduced from $N$ to $c\sqrt N$.
Iterating square roots takes only $O(\log\log N)$ steps, so the total growth is
at most a power of $\log N$.  The remaining general positive-operator
question is open because the coordinate-block test uses the shift geometry.

\section{The partial theorem}

\begin{theorem}\label{thm:main}
Let $1\le p\le\infty$, and let $T$ be the weighted backward shift on
$\ell^p(\N_0)$ defined by a bounded sequence $(w_j)$ satisfying $w_j\ge1$.
If $T$ is strongly Kreiss bounded, then there are constants $C,\kappa>0$
such that
\[
  \|T^N\|\le C\log^\kappa(N+2)\qquad(N\ge1).
\]
In particular, Question 5.8 has an affirmative answer in this class.
\end{theorem}

We first record the only scalar estimate needed.

\begin{lemma}[Poisson window]\label{lem:poisson}
There are constants $\eta,c_0>0$ and $N_0\in\N$ such that, whenever
$N\ge N_0$ and $m\in\N_0$ satisfies $|m-N|\le\eta\sqrt N$, one has
\[
  e^{-N}\frac{N^m}{m!}\ge \frac{c_0}{\sqrt N}.
\]
\end{lemma}

\begin{proof}
By the upper Stirling estimate, for an absolute $C>0$,
\[
 m!\le C\sqrt m\,(m/e)^m.
\]
Writing $m=N+s$, the exponential factor in the resulting lower bound is
\[
 -N+m+m\log(N/m)
 =s-(N+s)\log(1+s/N).
\]
For $|s|\le\eta\sqrt N$ this expression is bounded below by an absolute
constant, while $m$ is comparable with $N$.  Decreasing $c_0$ handles the
finitely many remaining values after fixing $N_0$.
\end{proof}

\begin{lemma}[Square-root recurrence]\label{lem:recurrence}
Under the assumptions of Theorem \ref{thm:main}, define
\[
  b_N:=\|T^N\|=\sup_{n\ge0}\frac{a_{n+N}}{a_n}.
\]
There are $C_0\ge1$, $c\in(0,1)$, and $N_1\in\N$ such that
\begin{equation}\label{eq:recurrence}
  b_N\le C_0 b_{\lfloor c\sqrt N\rfloor}\qquad(N\ge N_1).
\end{equation}
\end{lemma}

\begin{proof}
Let $\eta,c_0,N_0$ be as in Lemma \ref{lem:poisson}.  For large $N$, put
$L=\lfloor\eta\sqrt N/2\rfloor$.  From \eqref{eq:exp}, the normalized
Poisson operator
\[
  P_N:=e^{-N}e^{NT}
       =\sum_{m=0}^\infty e^{-N}\frac{N^m}{m!}T^m
\]
satisfies $\|P_N\|\le\K$.

Fix $n\ge0$.  If $p<\infty$, let
\[
 x=L^{-1/p}\sum_{k=n+N}^{n+N+L-1}e_k;
\]
if $p=\infty$, omit the factor $L^{-1/p}$.  In either case $\|x\|_p=1$.
For $j\in\{n,\ldots,n+L-1\}$ and every
$k\in\supp x$, the integer $m=k-j$ satisfies $|m-N|<L$, and hence its
Poisson coefficient is at least $c_0/\sqrt N$.  Since $w_j\ge1$, the sequence
$(a_j)$ is nondecreasing, and therefore
\[
 \frac{a_k}{a_j}\ge\frac{a_{n+N}}{a_{n+L}}.
\]
Using \eqref{eq:powers}, all terms are nonnegative and we obtain
\[
 (P_Nx)_j
 \ge \frac{c_0L}{\sqrt N}\,
      \frac{a_{n+N}}{a_{n+L}}L^{-1/p}.
\]
Taking the $\ell^p$ norm over the $L$ displayed output coordinates (with the
usual interpretation for $p=\infty$), and using $L/\sqrt N\ge\eta/4$ for
large $N$, gives
\[
 \K\ge\|P_Nx\|_p
 \ge \frac{c_0\eta}{4}\frac{a_{n+N}}{a_{n+L}}.
\]
Thus $a_{n+N}/a_{n+L}\le C_0:=4\K/(c_0\eta)$, uniformly in $n$.  Consequently,
\[
 \frac{a_{n+N}}{a_n}
 =\frac{a_{n+N}}{a_{n+L}}\frac{a_{n+L}}{a_n}
 \le C_0 b_L.
\]
Taking the supremum over $n$ proves \eqref{eq:recurrence}, after changing
$c$ and the finite threshold if necessary.
\end{proof}

\begin{proof}[Proof of Theorem \ref{thm:main}]
The numbers $b_N$ are nondecreasing because $(a_j)$ is nondecreasing.  Starting
from $N$, repeatedly apply Lemma \ref{lem:recurrence} while the current index
is at least $N_1$.  The indices obey $N_{r+1}\le c\sqrt{N_r}$, so the process
stops after at most
\[
  r\le C_1+\log_2\log(N+2)
\]
steps, where $C_1$ depends only on $c$ and $N_1$.  If
$B=\max_{1\le m\le N_1}b_m$, iteration yields
\[
 b_N\le B C_0^r
 \le C\,\log(N+2)^{\log_2 C_0}.
\]
This proves the theorem with $\kappa=\log_2 C_0$ (or any larger positive
number).
\end{proof}

\section{Verification and limitations}

The proof is independent of computation.  The included script
\texttt{code/weighted\_shift\_probe.py} computes finite-section induced norms
of $e^{-r}e^{rT}$ for cumulative weights $a_n=(n+1)^\beta$, three values of
$p$, three values of $\beta$, and four radii.  The observed norms increase
with the radius, as predicted by the theorem because polynomial power growth
cannot coexist with a uniform strong Kreiss bound in this monotone shift
class.  This finite computation is only a sanity check.

The result does not cover arbitrary positive operators on $L^p$, nor weighted
shifts with nonmonotone cumulative products.  A bounded literature search
through 19 July 2026 used the phrases ``positive strongly Kreiss bounded'',
``strongly Kreiss weighted shift'', and ``weighted shift logarithmic powers'',
and checked \cite{ArnoldCuny,CCEL,BermudezEtAl}.  No statement of Lemma
\ref{lem:recurrence} or Theorem \ref{thm:main} was found.  Novelty confidence
is moderate because the argument may be folklore in weighted-shift theory.

The recommended human-review focus is the uniform Poisson lower bound and the
two block endpoints in Lemma \ref{lem:recurrence}.  Once that lemma is
accepted, the logarithmic iteration is immediate.

\begin{thebibliography}{9}

\bibitem{ArnoldCuny}
L. Arnold and C. Cuny,
\emph{On the growth rate of powers of a strongly Kreiss bounded operator on
an $L^p$-space}, Studia Math., published online 23 March 2026,
DOI: \href{https://doi.org/10.4064/sm230303-29-10}{10.4064/sm230303-29-10};
arXiv:2302.14135.

\bibitem{CCEL}
G. Cohen, C. Cuny, T. Eisner, and M. Lin,
\emph{Resolvent conditions and growth of powers of operators},
J. Math. Anal. Appl. 487 (2020), 124035; arXiv:1912.10507.

\bibitem{BermudezEtAl}
T. Bermudez, A. Bonilla, V. Muller, and A. Peris,
\emph{Cesaro bounded operators in Banach spaces},
J. Math. Anal. Appl. 472 (2019), 1264--1278; arXiv:1706.03638.

\end{thebibliography}

\end{document}
